% !TeX root = ../main.tex

% ==================================================
% 2.8. 상속과 메서드 오버라이딩 — 자식 클래스는 부모 클래스의 기능을 어떻게 사용하는가?
% ==================================================

\subsection{상속과 메서드 오버라이딩 - 자식 클래스는 부모 클래스의 기능을 어떻게 사용하는가?}

2.7에서는 Python이 인스턴스 속성을 읽을 때 인스턴스 네임스페이스뿐 아니라 클래스 네임스페이스도 탐색한다는 사실을 학습했다.

이번 2.8에서는 이 속성 탐색 범위를 하나의 클래스에서 여러 클래스 사이의 관계로 확장한다.

다음 코드를 살펴보자.

\begin{dhkcode}
class Animal:
    def speak(self):
        print("Animal sound")


class Dog(Animal):
    pass

dog = Dog()
dog.speak()
\end{dhkcode}

실행 결과:

\begin{dhkzsh}
Animal sound
\end{dhkzsh}

여기서 중요한 질문이 생긴다.

\zsh{Dog} 클래스에는 \zsh{speak()}가 정의되어 있지 않다.

그런데 왜 \zsh{dog.speak()}를 호출할 수 있을까?

이번 학습에서는 이 질문을 출발점으로 다음 개념들을 연결할 것이다.

\begin{dhksmry}
    - 상속(Inheritance)과 클래스 간 관계

    - 부모 클래스와 자식 클래스의 객체 모델

    - 상속된 속성과 메서드의 탐색

    - 메서드 오버라이딩(Overriding)

    - self와 실제 인스턴스 타입

    - __init__()의 상속과 초기화

    - 다형성(Polymorphism)

    - PyTorch nn.Module의 상속 구조
\end{dhksmry}

\subsubsection{개념 정의 - 상속이란 무엇인가?}

상속(Inheritance)은 기존 클래스의 속성과 동작을 활용하여 새로운 클래스를 정의하는 메커니즘이다.

다음 코드를 보자.

\begin{dhkcode}
class Animal:
    def speak(self):
        print("Animal sound")


class Dog(Animal):
    pass
\end{dhkcode}

여기서 \zsh{Animal}은 부모 클래스(Base Class), \zsh{Dog}는 자식 클래스(Derived Class)이다.

\begin{dhkcode}
class Dog(Animal):
\end{dhkcode}

이 구문은 \zsh{Dog} 클래스가 \zsh{Animal}을 상속하도록 정의한다는 의미다.

상속 관계를 그림으로 표현하면 다음과 같다.

\begin{dhkgraph}[uniform]
node n1 (1,1) [accent] :: speak -> 함수 객체
node n2 (2,1) :: 부모의 speak() 사용 가능
group g1(n1) :: Animal 클래스 객체 
group g2(n2) :: Dog 클래스 객체 
edge g2 -> g1

caption :: 상속은 클래스 객체 사이에 속성 탐색 관계를 형성한다.
\end{dhkgraph}

여기서 가장 먼저 이해해야 할 점이 있다.

상속은 부모 클래스의 모든 속성과 메서드를 자식 클래스의 네임스페이스에 복사하는 과정이 아니다.

자식 클래스는 부모 클래스를 참조하는 상속 관계를 가지고, 속성을 탐색할 때 그 관계를 이용한다.

이 사실은 이후 메서드 오버라이딩과 \zsh{super()}를 이해하는 핵심이 된다.

\subsubsection{상속하면 실제로 무엇이 생성되는가?}

2.2에서는 \zsh{class}문을 실행하면 클래스 객체가 생성된다고 학습했다.

상속을 사용해도 마찬가지다.

\begin{dhkcode}
class Animal:
    pass

class Dog(Animal):
    pass
\end{dhkcode}

이 코드를 실행하면 두 개의 클래스 객체가 생성된다.

\begin{dhksmry}
    - Animal 클래스 객체

    - Dog 클래스 객체
\end{dhksmry}

\zsh{Dog}가 \zsh{Animal}을 상속했다고 해서 두 이름이 같은 객체를 참조하는 것은 아니다.

\begin{dhkcode}
print(Dog is Animal)
\end{dhkcode}

실행 결과:

\begin{dhkzsh}
False
\end{dhkzsh}

두 클래스는 별개의 객체다.

그렇다면 상속 관계는 어디에 기록될까?

Python 클래스 객체에는 직접적인 부모 클래스 정보를 확인할 수 있는 

\zsh{__bases__} 속성이 있다.

\begin{dhkcode}
print(Dog.__bases__)
\end{dhkcode}

실행 결과:

\begin{dhkzsh}
(<class '__main__.Animal'>,)
\end{dhkzsh}

즉, \zsh{Dog} 클래스 객체는 자신의 직접적인 부모 클래스가 \zsh{Animal}이라는 정보를 가지고 있다.

반대로 다음 코드도 확인해 보자.

\begin{dhkcode}
print(Animal.__bases__)
\end{dhkcode}

실행 결과:

\begin{dhkzsh}
(<class 'object'>,)
\end{dhkzsh}

Python 3의 일반적인 사용자 정의 클래스는 별도의 부모 클래스를 명시하지 않으면 기본적으로 \zsh{object}를 상속한다.

따라서 다음 두 클래스 정의는 기본적인 상속 관계에서 동일하다.

\begin{dhkcode}
class Animal:
    pass


class Animal(object):
    pass
\end{dhkcode}

전체 관계는 다음과 같다.

\begin{dhkdiagram}[in=10]
node n1 (1,1) :: object
node n2 (3,1) :: Animal
node n3 (5,1) :: Dog

edge n2 -> n1 :: 상속
edge n3 -> n2 :: 상속
\end{dhkdiagram}

이 그림의 화살표는 클래스 사이의 상속 관계를 나타낸다.

2.3에서 학습했던 클래스와 인스턴스의 타입 관계와는 구분해야 한다.

\begin{dhknote}
\begin{dhksmry}[centered]
상속 관계와 인스턴스 관계 비교
\end{dhksmry}
\end{dhknote}

\begin{dhkcode}
dog = Dog()
\end{dhkcode}

이때 두 종류의 관계가 동시에 존재한다.

\begin{dhkdiagram}[in=10]
node n1 (1,1) :: object
node n2 (3,1) :: Animal
node n3 (5,1) :: Dog
node n4 (7,1) :: dog
edge n4 -> n3 :: 인스턴스 관계
edge n3 -> n2 :: 상속
edge n2 -> n1 :: 상속
\end{dhkdiagram}

따라서 다음 결과가 성립한다.

\begin{dhkcode}
print(type(dog) is Dog)

print(isinstance(dog, Dog))
print(isinstance(dog, Animal))
print(isinstance(dog, object))
\end{dhkcode}

실행 결과:

\begin{dhkzsh}
True
True
True
True
\end{dhkzsh}

\zsh{dog}의 실제 타입은 \zsh{Dog}이지만, 상속 관계 때문에 \zsh{Animal}과 \zsh{object}의 인스턴스로도 인정되는 것이다.

\subsubsection{부모 클래스의 메서드는 자식 클래스에 복사되는가?}

이번 학습에서 중요한 내용 중 하나다.

다음 코드를 보자.

\begin{dhkcode}
class Animal:
    def speak(self):
        print("Animal sound")


class Dog(Animal):
    pass
\end{dhkcode}

이제 클래스 네임스페이스를 확인해 보자.

\begin{dhkcode}
print("speak" in Animal.__dict__)
print("speak" in Dog.__dict__)
\end{dhkcode}

실행 결과:

\begin{dhkzsh}
True
False
\end{dhkzsh}

\zsh{Animal}의 네임스페이스에는 \zsh{speak} 함수 객체가 존재한다.

하지만 \zsh{Dog}의 네임스페이스에는 존재하지 않는다.

그런데도 다음 코드는 실행된다.

\begin{dhkcode}
dog = Dog()

dog.speak()
\end{dhkcode}

실행 결과:

\begin{dhkzsh}
Animal sound
\end{dhkzsh}

왜 그럴까?

Python이 \zsh{Dog}에서 \zsh{speak}을 찾지 못하면 상속 관계에 따라 \zsh{Animal}에서도 탐색하기 때문이다.

\begin{dhkgraph}[uniform]
node n1 (1,1) :: dog.speak() // 인스턴스를 통한 속성 접근
node n2 (2,1) :: Dog 클래스 탐색 // speak 없음
node n3 (3,1) :: Animal 클래스 탐색 // speak 함수 객체 발견
node n4 (4,1) :: 바운드 메서드 호출 // self -> dog 인스턴스

edge n1 -> n2
edge n2 -> n3
edge n3 -> n4

caption :: 일반적인 상속 메서드 탐색 과정을 단순화한 그림
\end{dhkgraph}

여기서 2.5의 바운드 메서드 개념이 연결된다.

\zsh{Animal}에서 발견한 함수 객체를 \zsh{dog} 인스턴스에 바인딩하는 것이다.

실제로 확인해 보자.

\begin{dhkcode}
method = dog.speak

print(method.__self__ is dog)
print(method.__func__ is Animal.__dict__["speak"])
\end{dhkcode}

실행 결과:

\begin{dhkzsh}
True
True
\end{dhkzsh}

즉, 부모 클래스에 정의된 함수 객체가 자식 클래스 인스턴스에 바인딩될 수 있다.

이것이 상속 메서드 호출의 핵심 원리다.

\subsubsection{부모 메서드의 \zsh{self}는 어떤 인스턴스를 참조하는가?}

다음 코드를 살펴보자.

\begin{dhkcode}
class Animal:
    def identify(self):
        print(type(self))


class Dog(Animal):
    pass


dog = Dog()
dog.identify()
\end{dhkcode}

실행 결과:

\begin{dhkzsh}
<class '__main__.Dog'>
\end{dhkzsh}

여기서 주목해야 할 점이 있다.

\zsh{identify()}는 \zsh{Animal} 클래스에 정의되어 있다.

그런데 \zsh{self}의 타입은 \zsh{Animal}이 아니라 \zsh{Dog}다.

왜 그럴까?

\zsh{self}는 메서드가 정의된 클래스를 의미하는 것이 아니라, 메서드 호출에 바인딩된 실제 인스턴스를 참조하기 때문이다.

다음 코드를 보자.

\begin{dhkcode}
animal = Animal()
dog = Dog()

animal.identify()
dog.identify()
\end{dhkcode}

실행 결과:

\begin{dhkzsh}
<class '__main__.Animal'>
<class '__main__.Dog'>
\end{dhkzsh}

두 호출은 같은 함수 객체를 사용하지만, 전달되는 인스턴스가 다르다.

개념적으로 다음과 같다.

\begin{dhkcode}
Animal.identify(animal)

Animal.identify(dog)
\end{dhkcode}

따라서 함수 본문은 같더라도 \zsh{self}가 참조하는 객체에 따라 동작 결과가 달라질 수 있다.

\begin{dhknote}
\begin{dhksmry}[centered]
부모 메서드에서 자식 인스턴스의 속성에 접근할 수도 있다
\end{dhksmry}
\end{dhknote}

\begin{dhkcode}
class Animal:
    def introduce(self):
        print(self.name)


class Dog(Animal):
    def __init__(self, name):
        self.name = name


dog = Dog("Buddy")
dog.introduce()
\end{dhkcode}

실행 결과:

\begin{dhkzsh}
Buddy
\end{dhkzsh}

\zsh{introduce()}는 \zsh{Animal}에 정의되어 있지만, \zsh{self.name}은 \zsh{dog} 인스턴스의 속성을 읽는다.

즉, 메서드가 정의된 위치와 메서드가 사용하는 인스턴스의 상태는 구분되는 개념이다.

\subsubsection{메서드 오버라이딩 - 자식 클래스에서 메서드를 재정의하기}

이제 오버라이딩(Overriding)을 살펴보자.

오버라이딩은 자식 클래스가 부모 클래스에 존재하는 메서드와 같은 이름의 메서드를 정의하여, 자식 클래스에서의 동작을 변경하는 것이다.

다음 코드를 보자.

\begin{dhkcode}
class Animal:
    def speak(self):
        print("Animal sound")


class Dog(Animal):
    def speak(self):
        print("Woof!")
\end{dhkcode}

이제 인스턴스를 생성해 보자.

\begin{dhkcode}
animal = Animal()
dog = Dog()

animal.speak()
dog.speak()
\end{dhkcode}

실행 결과:

\begin{dhkzsh}
Animal sound
Woof!
\end{dhkzsh}

여기서 \zsh{Dog}는 부모 클래스의 \zsh{speak()}을 변경한 것일까?

정확히는 아니다.

\zsh{Animal} 클래스의 \zsh{speak} 함수 객체는 그대로 존재한다.

\zsh{Dog} 클래스의 네임스페이스에 새로운 \zsh{speak} 함수 객체가 정의된 것이다.

확인해 보자.

\begin{dhkcode}
print(Animal.__dict__["speak"])
print(Dog.__dict__["speak"])

print(
    Animal.__dict__["speak"] is Dog.__dict__["speak"]
)
\end{dhkcode}

실행 결과의 마지막 줄:

\begin{dhkzsh}
False
\end{dhkzsh}

두 함수 객체는 서로 다르다.

따라서 구조는 다음과 같다.

\begin{dhkgraph}
node n1 (1,1):: speak -> 함수 객체 A
node n2 (2,1) [accent] :: speak -> 함수 객체 B
group g1(n1) :: Animal 클래스
group g2(n2) :: Dog 클래스 
edge g2 -> g1

caption :: 오버라이딩은 부모의 함수 객체를 수정하지 않고 자식 클래스에 같은 이름의 속성을 정의한다.
\end{dhkgraph}

따라서 \zsh{dog.speak()}에서는 \zsh{Dog}클래스의 \zsh{speak}가 먼저 발견된다.

부모 클래스의 \zsh{speak}까지 탐색할 필요가 없는 것이다.

\begin{dhknote}
\begin{dhksmry}[centered]
오버라이딩은 이름 탐색의 결과다
\end{dhksmry}
\end{dhknote}

이전에 인스턴스 속성이 클래스 속성을 가리는 Shadowing을 학습했다.

메서드 오버라이딩도 넓은 의미에서는 비슷한 원리다.

\begin{dhksmry}
    - 인스턴스 속성이 클래스 속성을 가릴 수 있다.

    - 자식 클래스의 속성이 부모 클래스의 속성을 가릴 수 있다.
\end{dhksmry}

다만 오버라이딩은 주로 상속 관계에서 메서드를 재정의하는 것을 의미한다.

\subsubsection{부모 메서드와 자식 메서드를 직접 비교하기}

다음 코드를 보자.

\begin{dhkcode}
class Animal:
    def speak(self):
        print("Animal sound")


class Dog(Animal):
    def speak(self):
        print("Woof!")


dog = Dog()
\end{dhkcode}

일반적으로 다음 코드를 실행하면

\begin{dhkcode}
dog.speak()
\end{dhkcode}

\zsh{Dog.speak()}가 호출된다.

그런데 부모 클래스에 정의된 메서드를 직접 호출할 수도 있다.

\begin{dhkcode}
Animal.speak(dog)
\end{dhkcode}

실행 결과:

\begin{dhkzsh}
Animal sound
\end{dhkzsh}

왜 가능할까?

2.5에서 배운 것처럼 클래스에서 함수 객체에 직접 접근한 뒤 인스턴스를 명시적으로 전달할 수 있기 때문이다.

\begin{dhkcode}
Animal.speak(dog)
\end{dhkcode}

이 표현식은 \zsh{Animal} 클래스에 정의된 함수 객체를 가져와 \zsh{dog} 인스턴스를 첫 번째 인수로 전달한다.

따라서 자식 클래스에서 오버라이딩했더라도 부모 클래스의 함수 객체 자체는 여전히 호출할 수 있다.

하지만 실무에서 부모 메서드를 호출할 때는 일반적으로 \zsh{super()}를 사용한다.

\begin{dhkcode}
class Dog(Animal):
    def speak(self):
        super().speak()
        print("Woof!")
\end{dhkcode}

\zsh{super()}는 단순히 부모 클래스 이름을 대체하는 문법이 아니다.

MRO(Method Resolution Order)를 기반으로 다음 메서드를 찾는 메커니즘이다.

이 부분은 2.9에서 자세히 다룰 것이다.

현재는 오버라이딩이 부모 메서드를 삭제하거나 수정하는 것이 아니라는 점을 기억하면 충분하다.

\subsubsection{\zsh{__init__()}도 상속되고 오버라이딩되는가?}

2.6에서는 \zsh{__init__()}이 생성된 인스턴스를 초기화하는 특수 메서드라고 학습했다.

그렇다면 \zsh{__init__()}도 상속될까?

그렇다.

\zsh{__init__()} 역시 클래스에 정의된 메서드이므로 상속 관계의 영향을 받는다.

\begin{dhknote}
\begin{dhksmry}[centered]
부모의 __init__()을 그대로 사용하는 경우
\end{dhksmry}
\end{dhknote}

\begin{dhkcode}
class Animal:
    def __init__(self, name):
        self.name = name


class Dog(Animal):
    pass


dog = Dog("Buddy")

print(dog.name)
\end{dhkcode}

실행 결과:

\begin{dhkzsh}
Buddy
\end{dhkzsh}

\zsh{Dog}에는 \zsh{__init__()}이 정의되어 있지 않다.

하지만 인스턴스를 생성할 때 상속 관계를 통해 \zsh{Animal.__init__()}이 사용된다.

따라서 \zsh{dog} 인스턴스에 \zsh{name} 속성이 설정되는 것이다.

\begin{dhknote}
\begin{dhksmry}[centered]
자식이 __init__()을 오버라이딩하는 경우
\end{dhksmry}
\end{dhknote}

\begin{dhkcode}
class Animal:
    def __init__(self, name):
        self.name = name

class Dog(Animal):
    def __init__(self, name, breed):
        self.name = name
        self.breed = breed
\end{dhkcode}

이제 \zsh{Dog}에는 자체적인 \zsh{__init__()}이 있다.

\begin{dhkcode}
dog = Dog("Buddy", "Beagle")
\end{dhkcode}

실행 결과:

\begin{dhkzsh}
{'name': 'Buddy', 'breed': 'Beagle'}
\end{dhkzsh}

여기서는 자식 클래스가 직접 \zsh{name}과 \zsh{breed}를 초기화했다.

그런데 다음 코드에서는 문제가 발생할 수 있다.

\begin{dhkcode}
class Animal:
    def __init__(self, name):
        self.name = name
        self.alive = True


class Dog(Animal):
    def __init__(self, name, breed):
        self.breed = breed

dog = Dog("Buddy", "Beagle")

print(dog.alive)
\end{dhkcode}

실행 결과:

\begin{dhkzsh}
AttributeError: 'Dog' object has no attribute 'alive'
\end{dhkzsh}

왜 이런 문제가 발생할까?

\zsh{Dog.__init__()}이 \zsh{Animal.__init__()}을 오버라이딩했고, 부모의 초기화 메서드를 호출하지 않았기 때문이다.

자식이 \zsh{__init__()}을 정의하면 부모의 \zsh{__init__()}이 자동으로 추가 실행되는 것은 아니다.

이 점은 반드시 기억해야 한다.

\begin{dhknote}
\begin{dhksmry}[centered]
부모 초기화를 활용하는 경우
\end{dhksmry}
\end{dhknote}

\begin{dhkcode}
class Animal:
    def __init__(self, name):
        self.name = name
        self.alive = True


class Dog(Animal):
    def __init__(self, name, breed):
        super().__init__(name)
        self.breed = breed
\end{dhkcode}

이제 다음 코드가 정상적으로 실행된다.

\begin{dhkcode}
dog = Dog("Buddy", "Beagle")

print(dog.__dict__)
\end{dhkcode}

실행 결과:

\begin{dhkzsh}
{
    'name': 'Buddy',
    'alive': True,
    'breed': 'Beagle'
}
\end{dhkzsh}

\zsh{super().__init__(name)}을 통해 상속 계층의 다음 초기화 메서드를 호출했기 때문이다.

이 구조는 PyTorch의 \zsh{super().__init__()}을 이해하는 데 직접적으로 연결된다.

\subsubsection{다형성 - 같은 메서드 호출이 다른 동작을 수행하는 이유}

이제 상속과 오버라이딩을 결합하면 중요한 객체지향 개념이 등장한다.

바로 \zsh{다형성(Polymorphism)}이다.

다형성은 같은 인터페이스를 통해 서로 다른 객체가 각자의 동작을 수행할 수 있다는 개념이다.

다음 코드를 보자.

\begin{dhkcode}
class Animal:
    def speak(self):
        print("Animal sound")


class Dog(Animal):
    def speak(self):
        print("Woof!")


class Cat(Animal):
    def speak(self):
        print("Meow!")
\end{dhkcode}

이제 서로 다른 클래스의 인스턴스를 생성해 보자.

\begin{dhkcode}
animals = [
    Animal(),
    Dog(),
    Cat()
]
\end{dhkcode}

다음 반복문을 실행하면 어떻게 될까?

\begin{dhkcode}
for animal in animals:
    animal.speak()
\end{dhkcode}

실행 결과:

\begin{dhkzsh}
Animal sound
Woof!
Meow!
\end{dhkzsh}

반복문에서는 항상 같은 표현식을 사용했다.

\begin{dhkcode}
animal.speak()
\end{dhkcode}

하지만 실제 호출되는 메서드는 인스턴스의 타입에 따라 달라진다.

왜 그럴까?

각 인스턴스가 속한 클래스에서 \zsh{speak} 속성을 탐색하고, 해당 함수 객체를 인스턴스에 바인딩하기 때문이다.

\begin{dhktable}
|인스턴스|	실제 타입|	선택되는 메서드|
|:---|:---|
|Animal()|	Animal|	Animal.speak|
|Dog()|	Dog|	Dog.speak|
|Cat()|	Cat|	Cat.speak|
\end{dhktable}

여기서 중요한 사실이 있다.

반복문은 \zsh{Dog}인지 \zsh{Cat}인지 직접 검사하지 않는다.

\begin{dhkcode}
if isinstance(animal, Dog):
    ...
elif isinstance(animal, Cat):
    ...
\end{dhkcode}

이런 조건문 없이도 객체가 제공하는 \zsh{speak()}을 호출하면 각자의 동작이 실행된다.

이것이 다형성의 핵심적인 활용 방식이다.

\begin{dhknote}
\begin{dhksmry}[centered]
다형성이 반드시 상속을 요구하는 것은 아니다
\end{dhksmry}
\end{dhknote}

Python에서는 다음과 같은 코드도 가능하다.

\begin{dhkcode}
class Robot:
    def speak(self):
        print("Beep!")


class Person:
    def speak(self):
        print("Hello!")
\end{dhkcode}

두 클래스 사이에는 직접적인 상속 관계가 없다.

하지만 다음 함수는 두 객체를 모두 처리할 수 있다.

\begin{dhkcode}
def make_speak(obj):
    obj.speak()


make_speak(Robot())
make_speak(Person())
\end{dhkcode}

실행 결과:

\begin{dhkzsh}
Beep!
Hello!
\end{dhkzsh}

Python에서는 이런 방식을 흔히 덕 타이핑(Duck Typing)이라고 부른다.

객체의 구체적인 타입보다 필요한 동작을 제공하는가에 집중하는 방식이다.

따라서 다형성은 상속과 오버라이딩으로 구현할 수도 있지만, 반드시 상속을 통해서만 구현되는 것은 아니다.

\subsubsection{상속은 언제 사용하는 것이 적절한가?}

여기서 상속의 설계 목적을 생각해 보자.

상속은 단순히 중복 코드를 줄이기 위한 문법만은 아니다.

기존 클래스가 제공하는 인터페이스와 동작을 유지하거나 확장하면서, 새로운 타입을 정의하는 수단이다.

예를 들어 다음 관계는 자연스럽다.

\begin{dhkcode}
class Animal:
    pass


class Dog(Animal):
    pass
\end{dhkcode}

\zsh{Dog}는 \zsh{Animal}의 한 종류라고 볼 수 있기 때문이다.

이를 흔히 IS-A 관계라고 설명한다.

반면 다음과 같은 관계는 상속보다 객체를 속성으로 보유하는 구성이 자연스러울 수 있다.

\begin{dhkcode}
class Engine:
    pass


class Car:
    def __init__(self):
        self.engine = Engine()
\end{dhkcode}

자동차는 엔진의 한 종류가 아니다.

자동차가 엔진을 가지고 있는 관계이다.

이런 구성을 \zsh{합성(Composition)}이라고 한다.

\begin{dhktable}
|구분|	상속|	합성|
|:---|:---|:---|
|관계|	IS-A|	HAS-A|
|예시|	Dog(Animal)|	Car가 Engine을 보유|
|구현|	부모 클래스를 지정|	다른 객체를 속성으로 보유|
|주된 목적|	인터페이스와 동작의 확장·특수화|	여러 객체의 기능 조합|
\end{dhktable}

실제 프로그램 설계에서는 상속이 항상 합성보다 좋은 것도 아니고, 합성이 항상 상속보다 좋은 것도 아니다.

중요한 것은 두 클래스 사이에 어떤 관계를 표현하려는지를 먼저 결정하는 것이다.

이 구분은 PyTorch에서 특히 중요하다.

\subsubsection{PyTorch와의 연결 - \zsh{nn.Module} 상속과 계층의 합성}

이;제 다음 모델 코드를 살펴보자.

\begin{dhkcode}
import torch
import torch.nn as nn


class MyModel(nn.Module):
    def __init__(self):
        super().__init__()

        self.linear = nn.Linear(10, 5)

    def forward(self, x):
        return self.linear(x)
\end{dhkcode}

이 코드에는 상속과 합성이 동시에 사용되고 있다.

\begin{dhknote}
\begin{dhksmry}[centered]
상속 관계
\end{dhksmry}
\end{dhknote}

\begin{dhkcode}
class MyModel(nn.Module):
\end{dhkcode}

\zsh{MyModel}은 \zsh{nn.Module}을 상속한다.

따라서 \zsh{MyModel}은 \zsh{nn.Module}의 하위 클래스다.

\begin{dhkcode}
print(issubclass(MyModel, nn.Module))
\end{dhkcode}

실행 결과:

\begin{dhkzsh}
True
\end{dhkzsh}

그리고 생성된 모델 인스턴스 역시 \zsh{nn.Module}의 인스턴스로 인정된다.

\begin{dhkcode}
model = MyModel()

print(isinstance(model, nn.Module))
\end{dhkcode}

실행 결과:

\begin{dhkzsh}
True
\end{dhkzsh}

\zsh{MyModel}은 \zsh{nn.Module}의 기능을 활용할 수 있다.

예를 들어 다음과 같은 메서드들이 있다.

\begin{dhkcode}
model.parameters()
model.state_dict()
model.train()
model.eval()
model.to("cpu")
\end{dhkcode}

이러한 기능들은 \zsh{MyModel}에 직접 정의하지 않았어도 상속을 통해 사용할 수 있다.

\begin{dhknote}
\begin{dhksmry}[centered]
합성 관계
\end{dhksmry}
\end{dhknote}

다음 코드를 보자.

\begin{dhkcode}
self.linear = nn.Linear(10, 5)
\end{dhkcode}

여기서 \zsh{nn.Linear} 인스턴스를 생성해 모델의 하위 모듈로 등록하고 있다.

즉, \zsh{MyModel}이 \zsh{nn.Linear}를 상속하는 것이 아니다.

\zsh{MyModel} 인스턴스가 \zsh{nn.Linear} 인스턴스를 포함하는 구조다.

이를 그림으로 표현하면 다음과 같다.

\begin{dhkgraph}[uniform]
node n1 (1,1) [accent] :: nn.Module
node n2 (2,1) :: MyModel
node n3 (2,2) :: model
node n4 (3,2) [accent] :: nn.Linear

group g1(n1) :: 부모 클래스
group g2(n2) :: 자식 클래스
group g3(n3) :: MyModel 인스턴스
group g4(n4) :: 하위 모듈 인스턴스

edge g2 -> g1 :: 상속
edge g3 -> g2
edge g3 -> g4 :: 포함

caption :: Mymodel은 nn.Module을 상속하고, 생성된 모델 인스턴스는 nn.Linear 인스턴스를 하위 모듈로 포함한다.
\end{dhkgraph}

\begin{dhknote}
\begin{dhksmry}[centered]
forward()는 오버라이딩인가?
\end{dhksmry}
\end{dhknote}

\begin{dhkcode}
def forward(self, x):
    return self.linear(x)
\end{dhkcode}

그렇다.

\zsh{MyModel}은 \zsh{nn.Module}이 제공하는 \zsh{forward()} 인터페이스를 오버라이딩한다.

\zsh{nn.Module}의 기본 \zsh{forward()}는 사용자가 모델에 맞는 계산을 정의하도록 마련된 메서드다.

따라서 \zsh{MyModel}은 자신의 연산을 \zsh{forward()}에 정의한다.

\begin{dhkcode}
x = torch.randn(2, 10)

output = model(x)
\end{dhkcode}

이때 \zsh{model(x)}는 \zsh{nn.Module}의 호출 메커니즘을 사용하고, 그 과정에서 실제 모델 인스턴스에 정의된 \zsh{MyModel.forward()}가 호출된다.

이것은 상속된 호출 메커니즘과 오버라이딩된 메서드가 결합되는 사례다.

다만 \zsh{model(x)}와 \zsh{model.forward(x)}는 내부 처리 경로가 완전히 같지는 않다.

\zsh{model(x)}는 PyTorch의 Hook 등 추가적인 처리를 포함할 수 있기 때문이다.

\subsubsection{핵심 정리}

이번 학습의 핵심을 정리해 보자.

\begin{dhknote}
\begin{dhksmry}[centered]
상속은 클래스 객체 사이에 관계를 형성한다.
\end{dhksmry}
\end{dhknote}

\begin{dhkcode}
class Dog(Animal):
    pass
\end{dhkcode}

\zsh{Dog}는 \zsh{Animal}을 상속하는 별도의 클래스 객체다.

\begin{dhknote}
\begin{dhksmry}[centered]
부모 클래스의 메서드는 자식 클래스에 복사되지 않는다.
\end{dhksmry}
\end{dhknote}

자식 클래스는 상속 관계를 따라 부모 클래스의 속성과 메서드를 탐색할 수 있다.

\begin{dhknote}
\begin{dhksmry}[centered]
상속된 메서드의 self는 실제 호출 대상 인스턴스를 참조한다.
\end{dhksmry}
\end{dhknote}

\begin{dhkcode}
dog.speak()
\end{dhkcode}

부모 클래스에 정의된 메서드라도 \zsh{self}는 \zsh{dog}를 참조한다.

\begin{dhknote}
\begin{dhksmry}[centered]
오버라이딩은 

자식 클래스에 같은 이름의 메서드를 새로 정의하는 것이다.
\end{dhksmry}
\end{dhknote}

\begin{dhkcode}
class Dog(Animal):
    def speak(self):
        print("Woof!")
\end{dhkcode}

부모의 함수 객체를 수정하는 것은 아니다.

\begin{dhknote}
\begin{dhksmry}[centered]
__init__()도 상속과 오버라이딩의 영향을 받는다.
\end{dhksmry}
\end{dhknote}

자식 클래스에서 \zsh{__init__()}을 재정의하면 부모 초기화가 자동으로 추가 실행되는 것은 아니다.

\begin{dhknote}
\begin{dhksmry}[centered]
super()는 상속 계층의 다른 구현을 활용할 때 사용한다.
\end{dhksmry}
\end{dhknote}

\begin{dhkcode}
super().__init__()
\end{dhkcode}

정확한 탐색 방식은 MRO에 따라 결정된다.

\begin{dhknote}
\begin{dhksmry}[centered]
다형성은 동일한 인터페이스로 

서로 다른 객체의 동작을 호출할 수 있게 한다.
\end{dhksmry}
\end{dhknote}

\begin{dhkcode}
animal.speak()
\end{dhkcode}

실제 객체 타입에 따라 선택되는 메서드가 달라질 수 있다.

\begin{dhknote}
\begin{dhksmry}[centered]
상속과 합성은 서로 다른 객체 설계 방식이다.
\end{dhksmry}
\end{dhknote}

\begin{dhksmry}
    - 상속 : IS-A 관계

    - 합성 : HAS-A 관계
\end{dhksmry}

\begin{dhknote}
\begin{dhksmry}[centered]
PyTorch는 상속과 합성을 함께 사용한다.
\end{dhksmry}
\end{dhknote}

\begin{dhkcode}
class MyModel(nn.Module):
\end{dhkcode}

\zsh{nn.Module}을 상속하고,

\begin{dhkcode}
self.linear = nn.Linear(10, 5)
\end{dhkcode}

하위 모듈 인스턴스를 포함하며,

\begin{dhkcode}
def forward(self, x):
\end{dhkcode}

\zsh{forward()}를 오버라이딩한다.

\subsubsection{2.8에서 반드시 이해해야 할 결론}

이번 학습의 핵심 구조는 다음과 같다.

\begin{dhksmry}[centered]
Instance Attribute Access
\da
Current Class
\da
Inheritance Hierarchy
\da
Method Resolution
\da
Bound Method
\da
Execution
\end{dhksmry}

이것은 일반적인 인스턴스 메서드 탐색을 단순화한 흐름이다.

실제 속성 탐색에는 2.7에서 학습한 디스크립터 우선순위도 적용된다.

가장 중요한 사실은 다음 두 가지다.

\begin{dhksmry}
첫째, 상속은 부모 클래스의 메서드를 복사하는 것이 아니라,

부모 클래스의 속성을 탐색할 수 있게 관계를 만드는 것이다.

둘째, 오버라이딩은 부모 클래스의 메서드를 수정하는 것이 아니라,

자식 클래스의 네임스페이스에 같은 이름의 새로운 메서드를 정의하는 것이다.
\end{dhksmry}

이 두 원리를 이해하면 PyTorch에서 \zsh{nn.Module}을 상속하고 \zsh{forward()}를 오버라이딩하는 구조를 자연스럽게 해석할 수 있다.

다음 2.9에서는 지금까지 설명한 상속 계층의 메서드 탐색을 더 엄밀하게 다룰 것이다.

특히 다중 상속에서 메서드 탐색 순서가 어떻게 결정되는지,

\zsh{super()}가 왜 단순히 부모 클래스를 호출하는 문법이 아닌지,

그리고 \zsh{super().__init__()}이 어떤 규칙으로 다음 메서드를 선택하는지 살펴볼 것이다.